\section*{Abstract（要旨）}
\addcontentsline{toc}{section}{Abstract（要旨）}

那須地域は日本有数のリゾート観光地として発展してきた一方、近年では人口減少、高齢化、地域活動の縮小など、地方地域共通の衰退構造も抱えている。

そのような地域課題に対し、観光客の受け皿として地域を支える大型宿泊施設に着目し、大型宿泊施設・観光客・地域住民という三者の関係を、従来の「消費・労働・競合関係」ではなく、「相互創生構造」として再構築することを目的とした研究と、その実践結果をまとめたものである。

従来、観光施設と地域住民は、競争・消費・サービス提供の関係として語られることが多かった。しかし実際には、大型宿泊施設が観光基盤を支える一方、地域住民も自身の趣味・文化活動・地域交流を通じて観光体験を豊かにし、双方が相互に価値を生成する構造が存在している。

本研究では、ホテルサンバレー那須において実施された地域住民参加型クリスマス聖歌企画を事例として分析を行った。本企画では、地域住民による自主的な合唱活動をホテル施設内で実施し、さらに観光客自身も歌唱へ参加する構造を形成した。

特に本企画では、

\begin{itemize}
\item かつて結婚式場として利用されていた太陽の教会
\item パイプオルガンを含む本格的な教会空間
\item 高齢者を中心とした地域住民による合唱
\item 観光客も共に歌う参加型構造
\item 聖歌という曲の背景解説
\item 歌詞カードと共に、主催者メッセージによる感謝
\end{itemize}

などを通じて、単なる鑑賞型イベントではない、「共に時間を過ごす観光体験」が形成された。

また、太陽の教会は、豪華なオルガンや内外装を備えた象徴的空間である一方、少子化や結婚式需要の減少により、以前ほど活用されなくなっていた。ホテル側にとっても、この空間を新たな地域文化の場として再活用することは重要な課題であった。

さらに、クリスマス聖歌の背景説明やサンタクロースの由来などを短く挟むことで、観光客は単に音楽を聴くだけではなく、「なぜこの歌が長く受け継がれてきたのか」を感覚的に理解しながら参加する構造が生まれていた。

分析の結果、

\begin{itemize}
\item 地域住民による文化参与
\item 観光客との共同参加
\item 多世代交流
\item 文化の体験的継承
\item 地域と観光施設の境界緩和
\end{itemize}

が観察され、ホテルサンバレー那須が単なる宿泊施設ではなく、地域文化と観光客を接続する文化的媒介空間として機能していることが明らかとなった。

また本研究では、歌詞カード内に掲載された感謝のメッセージを通じ、「訪れる人」「迎える人」「暮らす人」が共に歌い、同じ時間を共有するという意味設計にも着目した。

これは単なるイベント運営ではなく、「情報設計による社会接続」として機能していたと考えられる。

本研究は、観光を「消費型体験」から「共同生成型体験」へ転換する、新たな地域共存モデルの可能性を提示するものである。
